\subsection{Ordinary common targets and a first coherent nonlinear repair}

\begin{proposition}[C19: \status{Proved}; ordinary pushouts can fail in low degree]
\label{sa08:C19}
For every odd prime $p$, over $k=\Z_{(p)}$ there are two marked
maps $G_p\to F$ between admissible standard theories whose ordinary
PD-DGA pushout fails eventual vertical exactness already in form
degree one. The failure persists after rationalization.
\end{proposition}
\begin{proof}
On cumulative coordinates use the two substitutions
\[
 z\longmapsto U^{[p]},\qquad
 z\longmapsto T^{[p]}-(T-V)^{[p]}.
\]
Each has endpoints $0,T^{[p]}$. Applying the same expression to
every cumulative coordinate gives the strict simplicial
PD-DGA maps by the marked free-module property. In simplex
degree one the rationalized pushout is the relative de Rham
algebra over $\Q[T]$ of
\[
 R=\Q[T,U,W]/(U^p+W^p-T^p),\qquad W=T-V.
\]
Rational divided powers are factorial powers, so the pushout
imposes exactly this degree-zero relation and its differential.

Put
\[
 L=U+W,\qquad
 Q=\sum_{j=0}^{p-1}(-1)^jU^{p-1-j}W^j,\qquad
 \theta=T^pQ\,dL.
\]
The identity $LQ=T^p$ implies $L\,dQ+Q\,dL=0$.
Wedge with $dL$ and multiply by $Q$ to get
$T^p dQ\wedge dL=0$. Hence $\theta$ is closed.
This calculation holds in the integral PD pushout as well:
the ordinary polynomial relation follows by multiplying
the PD relation by $p!$.

After inverting $T$, set $u=U/T$ and $w=W/T$. Then
$u^p+w^p=1$, the element $u+w$ is a unit, and
\[
 \theta=T^{2p}\,d\log(u+w).
\]
The logarithmic differential is not exact. To check this
directly, extend the coefficient field to contain the
$p$ distinct roots of $-1$ and compactify the affine
curve with coordinates $[U':W':Z]$, satisfying
$U'^p+W'^p=Z^p$. The coordinate $Z$ is a new
homogenizing coordinate, not the invertible coefficient
$T$. At an infinity point $[\beta:1:0]$ with
$\beta^p=-1$ and $\beta\ne-1$, put $W'=1$, $Z=t$.
The formal solution
$U'=\beta(1-t^p)^{1/p}$ gives
\[
 u+w=(U'+W')/Z=(\beta+1)t^{-1}+O(t^{p-1}).
\]
Thus $d\log(u+w)$ has residue $-1$. A derivative
of a Laurent series has residue zero, so the form
is not exact.

For every $a\geq0$, the integral closed form
$T^a\theta$ remains nonexact after this flat
localization and field extension. Its weight
is $2p+a$. The vertical $H^1$ of the pushout
is therefore nonzero in unbounded weights.
\end{proof}

This refutes automatic admissibility for arbitrary marked
choices in an ordinary common-target pushout. It does
not rule out specially chosen maps or a genuinely
derived PD pushout. In particular the defect is
rational in form degree one; increasing $p$ does
not move it into a high cochain degree.
The exact finite certificate
\texttt{B\_check\_fermat\_pushout.py}, with receipt
\texttt{B\_fermat\_pushout\_checks.json}, checks
$p=3,5,7$: factorization, differentiated identities,
the closedness syzygy, homogeneous weights, and
the formal curve expansion with residue $-1$.

\begin{proposition}[C20: \status{Proved}; an absolute quadratic repair and its next obstruction group]
\label{sa08:C20}
Let $M=K(k,r)$, $r\geq2$, and let
$x:M\to A^{w,1}$ be a simplicial $k$-linear map
into a termwise-flat eventual theory. Assume $w$
is large enough that $A^{w,1}$ is horizontally
acyclic, the cycle objects in weight $2w$ have
horizontal type $k[j]$ for $2\leq j\leq4$, and
the vertical maps onto those cycles are surjective.
These are finite eventual-range requirements.
Put $y=dx$.

The map from the free odd-to-even disk on $M$
extends simplicially over the following natural
repair. Add parameter modules
\[
 z:\Gamma^2M\text{ in form degree }3,\qquad
 u:M\otimes M\text{ in degree }2,\qquad
 v:\Lambda^2M\text{ in degree }1,
\]
all of weight $2w$, and impose the differential
\[
 dz=\Gamma^2(y),\qquad
 du=z\mu-xy,\qquad
 dv=u\,\mathrm{alt}+xx.
\]
Here $\mu:M\otimes M\to\Gamma^2M$ is multiplication,
$\mathrm{alt}(a\wedge b)=b\otimes a-a\otimes b$,
$xy(a\otimes b)=x(a)y(b)$, and
$xx(a\wedge b)=x(a)x(b)$.
The parameter-arity-two piece is contractible.

This first absolute repair does not by itself
iterate by connectivity. For
$P=\Gamma^2K(\Z,2)$, the possible next
degree-three odd-generator relation obstruction
group is
\[
 \operatorname{Hom}_{D(\Z)}
 \bigl(NP\otimes NP,\Z[7]\bigr)
 \cong(\Z/2)^2.
\]
No claim is made here that a particular actual
obstruction class in that group is nonzero.
\end{proposition}
\begin{proof}
In parameter arity two the uncorrected free
disk has the cochain complex
\[
 \Lambda^2M\xrightarrow{\mathrm{alt}}
 M\otimes M\xrightarrow{\mu}\Gamma^2M
\]
in form degrees $2,3,4$. The formulas above
give its identity cone, with alternating
signs, so it is contractible. They define
a differential: $d(xy)=\Gamma^2(y)\mu$,
$\mu\,\mathrm{alt}=0$, and
$d(xx)=xy\,\mathrm{alt}$.

We verify that the required primitives in
$A$ can be chosen as simplicial linear maps.
The first cycle map
$\Gamma^2M\to Z^4(A^{2w,*})$ factors through
$\Gamma^2(A^{w,1})$. The latter is
acyclic by Statement~\ref{sa08:C05}.
The normalized source is bounded below
with free terms, so this cycle map is
chain-nullhomotopic. Lift the nullhomotopy
through the degreewise surjection
$A^{2w,3}\to Z^4(A^{2w,*})$ to obtain $z$.

The resulting closed map
$M\otimes M\to Z^3(A^{2w,*})$ is nullhomotopic:
Eilenberg--Zilber identifies the source
in the derived category with $k[2r]$,
whereas the target has type $k[3]$ and
$2r>3$. The same lifting argument gives
$u$.

Finally, $\Lambda^2M$ has no normalized
terms in degrees at most $r$: the
standard Dold--Kan model has $M_n=0$
for $n<r$ and $M_r=k$. Its normalized
terms are free. Thus it has no derived
maps to the last cycle target $k[2]$,
because $r+1\geq3$. This gives $v$.
All extensions then follow from the free
PD algebra property.

The homotopy-lifting step just used is
elementary. If $f:P\to Z$ is chain-nullhomotopic,
$P$ has projective terms, and $B\to Z$ is
degreewise surjective, lift the components
of a nullhomotopy $h$ to $\widetilde h$ in
$B$. The map $d\widetilde h+\widetilde h d$
is a chain map lifting $f$. Normalization
is exact, and Dold--Kan returns the
simplicial linear map.

For the second assertion, use the
Dold--Kan basis of $K(\Z,r)_n$ indexed
by the $r$ cut positions of a monotone
surjection $[n]\to[r]$. A pair of such
basis vectors survives square-functor
normalization exactly when their cut
sets cover all $n$ positions.
For $r=2$, the normalized divided-square
complex has ranks $1,3,3$ in degrees
$2,3,4$, with differentials
\[
 d_3=(2,0,-2),\qquad
 d_4=
 \begin{pmatrix}
 0&-1&-1\\
 1&1&0\\
 0&-1&-1
 \end{pmatrix}.
\]
There are no other terms. The columns
of $d_4$ span the kernel of $d_3$,
and its kernel is generated by the
primitive vector $(-1,1,-1)$.
Consequently
\[
 H_*(NP)=\Z/2[2]\oplus\Z[4].
\]
Since $\Z$ has global dimension one,
this also gives its derived direct-sum
decomposition. Equivalently it follows
by elementary integral changes of
basis in the displayed complex.

The tensor square therefore has two
copies of $\Z/2$ in degree six, a
copy of $\Z/2$ in degrees four and
five, and $\Z$ in degree eight.
The two degree-six summands contribute
$\operatorname{Ext}^1_\Z(\Z/2,\Z)$
to maps into $\Z[7]$; every other
contribution is zero. This proves
the stated group.
\end{proof}

The first assertion concerns an absolute sphere
parameter. It does not solve the compatibility
problem for prescribed repairs on the boundary
of a horizontal cell. The second assertion
locates a possible obstruction, whose value
may depend on earlier primitive choices and
the coupled repairs of the full algebra.
Neither assertion supplies a complete
integral C.3 comparison.
Statement~\ref{sa08:C22} subsequently removes
the actual degree-three quadratic obstruction
by simultaneous choices, while retaining
the nonzero group calculated here.

The script
\texttt{B\_pd\_square\_normalization.py}
records the exact normalized matrices
for $1\leq r\leq4$, together with all chain
and vertical checks, in the JSON receipt of the same base name.
The next obstruction
calculation is independently reproduced from
the saved $r=2$ complex by
\texttt{B\_pd\_square\_next\_obstruction.py}.
Its JSON receipt includes the integral
matrices and unimodular Smith certificates;
the final nonunit factors are $2,2$.
